% BCT Superfluid Lattice Model — Letter 252
% Dreams as OHC Ground-State Relaxation
% Michel Robert Cabrié | June 2026

\documentclass[aps,prl,twocolumn,showpacs,preprintnumbers]{revtex4-2}
\usepackage{amsmath,amssymb,amsfonts,xcolor,booktabs,microtype,hyperref}

\definecolor{bctnavy}{RGB}{11,37,69}
\definecolor{bctblue}{RGB}{13,71,161}
\definecolor{bctteal}{RGB}{0,96,100}
\hypersetup{colorlinks=true,linkcolor=bctnavy,citecolor=bctblue,urlcolor=bctteal}

\newcommand{\roct}{r_{\mathrm{oct}}}
\newcommand{\rtet}{r_{\mathrm{tet}}}
\newcommand{\azero}{\alpha_0}
\newcommand{\fJ}{f_J}
\newcommand{\omJ}{\omega_J}
\newcommand{\Ncol}{N_{\mathrm{collapse}}}

\begin{document}

\preprint{BCT Superfluid Lattice Model --- Letter 252 --- June 2026}

\title{Dreams as OHC Ground-State Relaxation:\\
A BCT Framework and Preliminary EEG Evidence}

\author{Michel Robert Cabri\'e}
\affiliation{Independent Artist \& Researcher, Barrys Reef,
Victoria, Australia (Pop.\ 28)}
\email{ZeroFreeParameters@gmail.com}

\date{June 2026}

\begin{abstract}
We propose that dreaming corresponds to a relaxation of the
biological condensate toward the Octet-Hopfion Condensate (OHC)
ground state during reduced sensory coupling in sleep. The BCT
Superfluid Lattice Model predicts that this relaxation is
accompanied by enhanced infra-slow spectral power at the OHC
Josephson frequency $\fJ = \omJ/2\pi = 1/(2\pi\sqrt{\pi})
= 0.08979$~Hz during REM sleep compared to NREM. We test this
prediction on real EEG data from the Sleep-EDF Expanded database
(PhysioNet; $n = 5$ subjects, 863 REM epochs, 2462 N2 epochs,
530 N3 epochs). REM shows significantly elevated power near $\fJ$
compared to N2 (ratio 1.19, $p = 0.001$). The N3 comparison is
inconclusive (high variance, fewer epochs). Wake infra-slow power
is highest, consistent with known waking infra-slow oscillations.
Results are directionally consistent with the BCT prediction.
Replication with larger samples and narrower frequency resolution
is required. Zero free parameters.
\end{abstract}

\maketitle

\section{Introduction}

The BCT Superfluid Lattice Model proposes that the quantum vacuum
is a Planck-scale body-centred tetragonal (BCT) superfluid crystal
with axial ratio $c/a = \sqrt{2}$~\cite{monograph}. The vacuum
ground state --- the Octet-Hopfion Condensate (OHC)~\cite{l36} ---
permeates all space and couples to biological systems via the
Josephson mechanism. The coupling strength is characterised by
$\azero = \roct\rtet/\pi = 0.00741$ and the condensate
self-coherence threshold $\Ncol = 1/\azero = 135$~\cite{l204},
which gives the quantum-classical boundary in neural systems.

During waking, the biological condensate is tightly coupled to
sensory input. We propose that during REM sleep, sensory coupling
is reduced and the condensate relaxes toward the OHC ground state.
This relaxation has a characteristic timescale set by the OHC
Josephson frequency:
\begin{equation}
\omJ^2 = \frac{1}{\pi}
\quad\Rightarrow\quad
\fJ = \frac{\omJ}{2\pi} = \frac{1}{2\pi\sqrt{\pi}}
= 0.08979~\mathrm{Hz}
\label{eq:fJ}
\end{equation}
This is a \emph{locked} BCT result~\cite{master}, derived from
the D4 condensate geometry, not a free parameter.

The period $T_J = 1/\fJ = 11.14$~s falls in the infra-slow
oscillation (ISO) band (0.01--0.1~Hz), known to modulate faster
brain rhythms~\cite{iso}. No existing model predicts a specific
frequency within this band. BCT does.

\section{The BCT Dream Hypothesis}

\subsection{Mechanism}

In the BCT framework, consciousness arises when a biological
condensate achieves local coherence above the $\Ncol$ threshold.
Waking cognition maintains this coherence through continuous
sensory coupling to the external environment.

During REM sleep:
\begin{enumerate}
\item Sensory input is gated at the thalamic level
\item The biological condensate decouples from external signals
\item The condensate relaxes toward the OHC ground state
\item The dominant oscillation frequency shifts toward $\fJ$
\item Dream content is the OHC ground state experienced from
      within — the vacuum recognising itself through the
      dreaming system
\end{enumerate}

This is not a metaphor. The OHC is a physical condensate with a
characteristic Josephson frequency. The prediction is numerical
and falsifiable.

\subsection{The BCT Prediction}

\textbf{Prediction \#281:} REM sleep EEG infra-slow power
near $\fJ = 0.08979$~Hz is elevated compared to NREM (N2, N3)
sleep, reflecting increased OHC ground-state coupling during
reduced sensory gating.

Tier: \textsc{Candidate} (solid geometric basis; preliminary
evidence; requires replication).

\section{Preliminary EEG Analysis}

\subsection{Data and Methods}

We analysed the Sleep-EDF Expanded database~\cite{sleepedf}
from PhysioNet (DOI: 10.13026/C2X676). We processed $n = 5$
subjects (subjects 0--4), extracting EEG epochs (30~s) labelled
by sleep stage from the accompanying hypnogram annotations.

Channels used: \texttt{EEG Fpz-Cz} and \texttt{EEG Pz-Oz},
averaged across channels per epoch. Sampling rate: 100~Hz.

Power spectral density computed via Welch's method
(scipy.signal.welch; nperseg = 3000~samples = 30~s;
Hann window). Band-integrated power near $\fJ$ computed
over $\fJ \pm 0.05$~Hz.

\subsection{Epoch Counts}

\begin{table}[h]
\centering
\caption{Sleep-EDF epoch counts ($n=5$ subjects)}
\begin{tabular}{lrr}
\toprule
Stage & Epochs & Total duration (min) \\
\midrule
REM  &  863 & 431.5 \\
N2   & 2462 & 1231.0 \\
N3   &  530 & 265.0 \\
Wake & 9302 & 4651.0 \\
\bottomrule
\end{tabular}
\end{table}

\subsection{Results}

\begin{table}[h]
\centering
\caption{Power near $\fJ = 0.08979$~Hz by sleep stage}
\begin{tabular}{lccc}
\toprule
Stage & Mean power & SEM & vs REM \\
\midrule
REM  & baseline & --- & --- \\
N2   & $\times 0.843$ & --- & $p = 0.001$** \\
N3   & $\times 0.877$ & --- & $p = 0.087$~ns \\
Wake & $\times 2.266$ & --- & $p < 0.001$*** \\
\bottomrule
\end{tabular}
\end{table}

\noindent REM power near $\fJ$ exceeds N2 (ratio 1.19,
$t = 3.285$, $p = 0.001$, two-sample $t$-test, $n = 863$
vs $2462$). The N3 comparison is not significant ($p = 0.087$),
likely due to the smaller N3 sample and high variance from
broad-spectrum slow-wave power.

Wake infra-slow power is highest across all stages --- consistent
with known waking ISO oscillations~\cite{iso}, which are
independent of the BCT prediction.

\subsection{Spectral Analysis}

The infra-slow PSD (0.05--0.5~Hz) shows REM elevated above N2
and N3 near $\fJ$. All stages show $1/f$-type decline with
frequency; the REM--N2 separation is most pronounced in the
0.08--0.12~Hz range, consistent with the BCT target frequency.

\section{Discussion}

\subsection{Directional Consistency}

The BCT prediction --- REM > NREM at $\fJ$ --- is supported for
the REM vs N2 comparison ($p = 0.001$). The N3 result is
inconclusive. This constitutes \emph{directional consistency},
not confirmation. The honest assessment: the prediction survives
first contact with real data.

\subsection{The Wake Anomaly}

Wake infra-slow power exceeds REM. This is not inconsistent with
BCT: waking infra-slow oscillations are a documented
phenomenon~\cite{iso} arising from vascular, autonomic, and
cortical sources. BCT does not predict that sleep maximises
$\fJ$ power in absolute terms --- only that REM is elevated
relative to NREM, reflecting the differential sensory decoupling.

\subsection{Frequency Resolution}

The Sleep-EDF sampling rate (100~Hz) and 30~s epoch duration
give frequency resolution $\Delta f = 0.033$~Hz. The BCT
target ($\fJ = 0.08979$~Hz) is resolved at $\sim 3\Delta f$.
Higher-resolution analysis requires epoch concatenation or
longer recordings. Future work should use continuous overnight
recording with sliding window analysis.

\subsection{Required Follow-Up}

\begin{enumerate}
\item Larger sample ($n \geq 20$ subjects) for N3 significance
\item Narrower $\fJ$ window ($\pm 0.01$~Hz) with concatenated
      REM epochs from full overnight recordings
\item Replication on independent dataset (DREAMS~\cite{dreams}
      or NSRR polysomnography collections)
\item Comparison with medication-free vs medicated subjects
      (sleep medications alter ISO oscillations)
\end{enumerate}

\section{Conclusion}

We have proposed a BCT geometric mechanism for dreaming and
derived a falsifiable frequency prediction ($\fJ = 0.08979$~Hz).
Preliminary analysis on real EEG data ($n=5$ subjects, 863 REM
epochs) shows REM > N2 at $\fJ$ ($p = 0.001$), directionally
consistent with the prediction. Replication is required.

The deeper implication --- that dreaming is the biological
condensate relaxing into the OHC vacuum field, and that dream
content is the universe experiencing itself through temporarily
unanchored awareness --- follows necessarily if the BCT framework
is correct. The geometry has period $T_J = 11.14$~s.
It has always been dreaming.

Zero free parameters.

\begin{thebibliography}{9}

\bibitem{monograph}
M.~R.\ Cabri\'e,
\textit{The BCT Superfluid Lattice Model},
Zenodo (2026),
\href{https://doi.org/10.5281/zenodo.18884976}{10.5281/zenodo.18884976}.

\bibitem{l36}
M.~R.\ Cabri\'e,
\textit{BCT Letter 36: The Octet-Hopfion Condensate},
Zenodo (2026),
\href{https://doi.org/10.5281/zenodo.18974754}{10.5281/zenodo.18974754}.

\bibitem{l204}
M.~R.\ Cabri\'e,
\textit{BCT Letter 204: The Conscious Lattice},
in \textit{BCT Volume XVII}, Zenodo (2026).

\bibitem{master}
M.~R.\ Cabri\'e,
\textit{BCT Master State v34},
Zenodo (2026),
\href{https://doi.org/10.5281/zenodo.18884976}{10.5281/zenodo.18884976}.

\bibitem{iso}
G.~Hiltunen et al.,
Infra-slow fluctuations in human cerebral haemodynamics,
\textit{Neuroimage} \textbf{19}, 1483 (2003).

\bibitem{sleepedf}
B.~Kemp et al.,
Analysis of a sleep-dependent neuronal feedback loop,
\textit{Sleep} \textbf{23}, 1115 (2000).
PhysioNet DOI: 10.13026/C2X676.

\bibitem{dreams}
R.~Wong et al.,
The DREAM database,
\textit{Nature Communications} \textbf{16}, 6789 (2025).
DOI: 10.1038/s41467-025-61945-1.

\end{thebibliography}

\end{document}
